\documentclass[12pt,a4paper]{article}
%%%%%%%%%%%%%%%%%%%%%%%%%%%%%%%%%%%%%%%%%%%%
%                 PACKAGES
%%%%%%%%%%%%%%%%%%%%%%%%%%%%%%%%%%%%%%%%%%%%

% LANGUAGE AND ENCODING
\usepackage[utf8]{inputenc}      % Direct UTF-8 character input
\usepackage[english]{babel}      % Language-specific hyphenation
\usepackage{xstring}

% CORE MATHEMATICS
\usepackage{amsmath}             % AMS math environments (align, gather, etc.)
\usepackage{amssymb}             % AMS math symbols
\usepackage{amsthm}              % Theorem environments
\usepackage{amsfonts}            % AMS fonts
\usepackage{mathtools}           % Extensions to amsmath
\usepackage{esint}               % Extended integral symbols
\usepackage{cancel}              % Diagonal cancellation in math
\usepackage{tcolorbox}           % Good boxes
\tcbuselibrary{most}

\usepackage{lmodern}
\usepackage{bm}                  % Bold math symbols (Greek, etc.)
\usepackage{mathrsfs}            % Ralph Smith's script font
\usepackage{tensor}              % Tensor index notation

% UNITS AND QUANTITIES
\usepackage{siunitx}             % Professional number and unit formatting
\NewCommandCopy\uqty\qty         % Save siunitx's \qty as \uqty
\usepackage{silence}
\WarningFilter{siunitx}{Detected the "physics" package}
\ExplSyntaxOn
\msg_redirect_name:nnn { siunitx } { physics-pkg } { none }
\ExplSyntaxOff

% TABLES AND BOXES
\usepackage{tabulary}            % Tables with automatic width adjustment
\usepackage{mdframed}            % Framed boxes with backgrounds
% \usepackage{chemformula}       % Alternative chemistry notation (commented)

% MATH TOOLS
\usepackage{centernot}           % Centered negation symbol
\usepackage{dcolumn}             % Decimal alignment in tables
\usepackage[shortlabels]{enumitem} % Customizable lists

% GRAPHICS AND FIGURES
\usepackage{graphicx}            % Include images
\usepackage{float}               % Force figure placement with [H]
\usepackage{tikz}                % Powerful diagram drawing
\usetikzlibrary{positioning,arrows,trees,3d,angles} % TikZ libraries
\usepackage{pgfplots}            % Scientific plots based on tikz
\pgfplotsset{compat=1.18}        % PGFPlots compatibility mode
\usepackage{subcaption}          % Subfigures (a), (b), etc.
\usepackage{caption}             % Customize caption formatting
\usepackage{booktabs}            % Professional tables with \toprule

% MULTI-COLUMN AND ARRAYS
\usepackage{multicol}            % Multi-column text layout
\usepackage{array}               % Enhanced table column types

% CHEMISTRY
\usepackage{chemfig}             % Draw molecular structures
\usepackage[version=4]{mhchem}   % Chemical equations in text (\ce{H2O})

% CIRCUITS
\usepackage{circuitikz}          % Circuit diagrams with tikz

% TEXT FORMATTING
\usepackage{blindtext}           % Lorem ipsum placeholder text
\usepackage{titlesec}            % Customize section headings
\usepackage{fancyhdr}            % Custom headers and footers
\usepackage{csquotes}            % Context-sensitive quotes

% HYPERLINKS
\usepackage[colorlinks=true, allcolors=blue]{hyperref} % Clickable links

% PAGE GEOMETRY
\usepackage[top=1.0in,bottom=1.0in,left=1.0in,right=1.0in,marginparwidth=1.75cm]{geometry}

% CODE LISTINGS
\usepackage{xcolor}              % Color definitions
\usepackage{listings}            % Code syntax highlighting
\usepackage{xparse}              % Robust command definitions
\usepackage{braket}              % Dirac bra-ket notation

% PHYSICS PACKAGES (LOADED LAST)
\usepackage[italicdiff]{physics} % Original physics package (d in italics)
\usepackage{physics2}            % Modern physics package
\usephysicsmodule{ab}            % Auto-braces module
\usephysicsmodule{ab.legacy}     % Legacy \abs, \norm, \eval
\usephysicsmodule{nabla.legacy}  % \grad, \div, \curl operators
\usephysicsmodule{op.legacy}     % Operators like \tr, \Re, \Im
\usephysicsmodule{diagmat, xmat} % Diagonal and special matrices
\usephysicsmodule{braket, ab.braket} % Dirac notation
\let\lpc\laplacian               % Alias for Laplacian

%%%%%%%%%%%%%%%%%%%%%%%%%%%%%%%%%%%%%%%%%%%%
%             NEW COMMANDS
%%%%%%%%%%%%%%%%%%%%%%%%%%%%%%%%%%%%%%%%%%%%

%%% MATH OPERATORS
\DeclareMathOperator{\sinc}{sinc}
\DeclareMathOperator{\sgn}{sgn}

%%% INVERSE TRIGONOMETRIC OPERATORS
\renewcommand{\asin}{\sin^{-1}}
\renewcommand{\acos}{\cos^{-1}}
\renewcommand{\atan}{\tan^{-1}}
\renewcommand{\acsc}{\csc^{-1}}
\renewcommand{\asec}{\sec^{-1}}
\renewcommand{\acot}{\cot^{-1}}

\newcommand{\asinh}{\sinh^{-1}}
\newcommand{\acosh}{\cosh^{-1}}
\newcommand{\atanh}{\tanh^{-1}}
\newcommand{\acsch}{\csch^{-1}}
\newcommand{\asech}{\sech^{-1}}
\newcommand{\acoth}{\coth^{-1}}
% \DeclareMathOperator{\arcsin}{arcsin}
% \DeclareMathOperator{\arccos}{arccos}
% \DeclareMathOperator{\arctan}{arctan}
% \DeclareMathOperator{\arccsc}{arccsc}
% \DeclareMathOperator{\arcsec}{arcsec}
% \DeclareMathOperator{\arccot}{arccot}
\DeclareMathOperator{\arsinh}{arsinh}
\DeclareMathOperator{\arcosh}{arcosh}
\DeclareMathOperator{\artanh}{artanh}
\DeclareMathOperator{\arcsch}{arcsch}
\DeclareMathOperator{\arsech}{arsech}
\DeclareMathOperator{\arcoth}{arcoth}

%%% NUMBER SETS
\newcommand{\R}{\mathbb{R}}          % Real numbers
\newcommand{\Z}{\mathbb{Z}}          % Integers
\newcommand{\Q}{\mathbb{Q}}          % Rational numbers
\newcommand{\C}{\mathbb{C}}          % Complex numbers
\newcommand{\N}{\mathbb{N}}          % Natural numbers
\newcommand{\K}{\mathbb{K}}          % Field R or C
\newcommand{\bb}[1]{\mathbb{#1}}     % Generic blackboard bold

%%% PARTIAL DIFFERENTIALS
\newcommand{\p}{\partial}            % Partial symbol

%%% VECTOR CALCULUS
\newcommand{\krondelta}[1]{\delta_{#1}}
\newcommand{\levicivita}[1]{\varepsilon_{#1}}
\newcommand{\ihat}{\vu*{\imath}}
\newcommand{\jhat}{\vu*{\jmath}}
\newcommand{\khat}{\vu*{k}}
\newcommand{\rhohat}{\vu*{\rho}}
\newcommand{\phihat}{\vu*{\phi}}
\newcommand{\thetahat}{\vu*{\theta}}
\newcommand{\rhat}{\vu*{r}}
\newcommand{\inv}{^{-1}}
\newcommand{\vell}{\boldsymbol{\ell}}

%%% DIFFERENTIALS AND TRANSFORMS
\newcommand{\laplace}[1]{\mathscr{L}\qty{#1}}
\newcommand{\invlaplace}[1]{\mathscr{L}^{-1}\qty{#1}}
\newcommand{\unitstep}[1]{\mathscr{U}\qty{#1}}
\newcommand{\intfty}{\int_{-\infty}^{+\infty}}

%%% PARENTHESES
\newcommand{\floor}[1]{\lfloor #1 \rfloor}
\newcommand{\ceiling}[1]{\lceil #1 \rceil}
\newcommand{\round}[1]{\lfloor #1 \rceil}
\newcommand{\nvec}[3]{\left\langle #1 , #2 , #3 \right\rangle}
\newcommand{\tline}[1]{\hrule \vspace{#1} \hrule}
\newcommand{\blank}[1]{\hspace*{#1}}

%%% TRANSFORM OPERATORS
\newcommand{\Laplace}{\mathscr{L}}
\newcommand{\Fourier}{\mathscr{F}}

%%% LOGICAL OPERATORS
\newcommand{\union}{\cup}
\newcommand{\intersect}{\cap}
\newcommand{\defined}{:=}

%%% COMPLEX ANALYSIS
\newcommand{\Arg}{\operatorname{Arg}}
\newcommand{\Hol}{\operatorname{Hol}}
\newcommand{\Log}{\operatorname{Log}}

%%% LINEAR ALGEBRA
\newcommand{\Span}[1]{\operatorname{span}\qty(#1)}
\newcommand{\Rank}[1]{\operatorname{rank}\qty(#1)}
\newcommand{\nullity}[1]{\operatorname{nullity}\qty(#1)}
\DeclareMathOperator{\image}{im}

%%% TOPOLOGY
\newcommand{\Id}{\operatorname{Id}}
\newcommand{\Image}[1]{\operatorname{Im}\qty(#1)}
\newcommand{\comp}{^{c}}
\newcommand{\Prob}[1]{\operatorname{Pr}\qty(#1)}
\newcommand{\du}{\mathrm{d}}
\newcommand{\topo}[1]{\mathcal{#1}}

%%% MISCELLANEOUS
\newcommand{\notimplies}{\centernot\implies}

\NewDocumentCommand{\dimu}{O{\mkern-3mu} m}{
\left[
#1
\left[#2\right]
#1
\right]
}

\newcommand{\id}{\operatorname{id}}
\newcommand{\contradiction}{\Rightarrow\!\Leftarrow}
\newcommand{\transition}[3]{{_{#2}\mathbf{#1}_{#3}}}
\newcommand{\barline}{\vspace{10pt}\hrule\vspace{10pt}}
\newcommand{\eqques}{\stackrel{?}{=}}
\newcommand{\Gammaf}[1]{\Gamma\qty(#1)}
\renewcommand\qedsymbol{$\square$}


%%%%%%%%%%%%%%%%%%%%%%%%%%%%%%%%%%%%%%%%%%%%
%       CUSTOM HEADER / SECTION STYLES
%%%%%%%%%%%%%%%%%%%%%%%%%%%%%%%%%%%%%%%%%%%%

\fancypagestyle{titleheader}{
  \fancyhead{}
  \fancyhead{\vspace{2cm}\tline{3pt}\fancyfoot[c]{\thepage}}
  \renewcommand{\headrulewidth}{0pt}
}
\newcommand{\titlebars}{
  \maketitle
  \tline{3pt}
  \thispagestyle{titleheader}
}

%%% CUSTOM SUBSUBSUBSECTION (4-LEVEL HEADINGS)
\titleclass{\subsubsubsection}{straight}[\subsection]
\newcounter{subsubsubsection}[subsubsection]
\renewcommand\thesubsubsubsection{\thesubsubsection.\arabic{subsubsubsection}}
\renewcommand\theparagraph{\thesubsubsubsection.\arabic{paragraph}}
\titleformat{\subsubsubsection}{\normalfont\normalsize\bfseries}{\thesubsubsubsection}{1em}{}
\titlespacing*{\subsubsubsection}{0pt}{3.25ex plus 1ex minus .2ex}{1.5ex plus .2ex}

\makeatletter
\renewcommand\paragraph{\@startsection{paragraph}{5}{\z@}{3.25ex \@plus1ex \@minus.2ex}{-1em}{\normalfont\normalsize\bfseries}}
\renewcommand\subparagraph{\@startsection{subparagraph}{6}{\parindent}{3.25ex \@plus1ex \@minus .2ex}{-1em}{\normalfont\normalsize\bfseries}}
\def\toclevel@subsubsubsection{4}
\def\toclevel@paragraph{5}
\def\toclevel@subparagraph{6}
\def\l@subsubsubsection{\@dottedtocline{4}{7em}{4em}}
\def\l@paragraph{\@dottedtocline{5}{10em}{5em}}
\def\l@subparagraph{\@dottedtocline{6}{14em}{6em}}
\makeatother
\setcounter{secnumdepth}{4}
\setcounter{tocdepth}{4}

%%% CUSTOM SECTION TITLE FONT SIZES
% \titleformat*{\section}{\LARGE\bfseries}
% \titleformat*{\subsection}{\Large\bfseries}
% \titleformat*{\subsubsection}{\large\bfseries}
% \setlength{\parskip}{10pt}

%%%%%%%%%%%%%%%%%%%%%%%%%%%%%%%%%%%%%%%%%%%%
%             ENVIRONMENTS
%%%%%%%%%%%%%%%%%%%%%%%%%%%%%%%%%%%%%%%%%%%%

\newenvironment{amatrix}[1]{%
  \left(
  \begin{array}{@{}*{#1}{c}|c@{}}}{
\end{array}\right)}

\newenvironment{abmatrix}[1]{%
  \left[
  \begin{array}{@{}*{#1}{c}|c@{}}}{
\end{array}\right]}


%%%%%%%%%%%%%%%%%%%%%%%%%%%%%%%%%%%%%%%%%%%%
%          THEOREM STYLES
%%%%%%%%%%%%%%%%%%%%%%%%%%%%%%%%%%%%%%%%%%%%

% \theoremstyle{definition}
% \newtheorem{theorem}{Theorem}[subsection]
% \newtheorem{corollary}{Corollary}[theorem]
% \newtheorem{lemma}[theorem]{Lemma}
% \newtheorem{definition}{Definition}[subsection]
% \newtheorem{example}{Example}[subsection]

% \theoremstyle{remark}
% \newtheorem*{remark}{Remark}
% \newtheorem*{solution}{Solution}
% \newtheorem*{note}{Note}

%%%%%%%%%%%%%%%%%%%%%%%%%%%%%%%%%%%%%%%%%%%%
%        CODE LISTING STYLE
%%%%%%%%%%%%%%%%%%%%%%%%%%%%%%%%%%%%%%%%%%%%


% ============================================
% PYTHON STYLE
% ============================================
\definecolor{vscodeBackground}{RGB}{255,255,255}
\definecolor{vscodeForeground}{RGB}{10,10,10}
\definecolor{vscodeKeyword}{RGB}{60,60,200}
\definecolor{vscodeString}{RGB}{230,0,0}
\definecolor{vscodeComment}{RGB}{0,100,0}
\definecolor{vscodeFunction}{RGB}{220,220,170}
\definecolor{vscodeOperator}{RGB}{212,212,212}
\definecolor{vscodeNumber}{RGB}{181,206,168}

\lstdefinestyle{vscodepython}{
    backgroundcolor=\color{vscodeBackground},
    basicstyle=\ttfamily\footnotesize\color{vscodeForeground},
    commentstyle=\color{vscodeComment},
    keywordstyle=\color{vscodeKeyword}\bfseries,
    stringstyle=\color{vscodeString},
    numberstyle=\tiny\color{vscodeForeground},
    identifierstyle=\color{vscodeForeground},
    morekeywords={True, False, None, and, as, assert, async, await, break, class,
        continue, def, del, elif, else, except, finally, for, from, global, if,
        import, in, is, lambda, nonlocal, not, or, pass, raise, return, try,
        while, with, yield},
    frame=single,
    rulecolor=\color{vscodeForeground},
    breaklines=true,
    numbers=left,
    numbersep=5pt,
    showstringspaces=false,
    tabsize=4,
    captionpos=b,
}

% ============================================
% C/C++ STYLE (FIXED)
% ============================================
\definecolor{vscodeType}{RGB}{0,150,200}
\definecolor{vscodePreprocessor}{RGB}{150,0,150}
\definecolor{vscodeFunction}{RGB}{180,100,0}
\definecolor{vscodeOperator}{RGB}{180,60,180}

\lstdefinestyle{vscodecpp}{
    backgroundcolor=\color{vscodeBackground},
    basicstyle=\ttfamily\scriptsize\color{vscodeForeground},
    commentstyle=\color{vscodeComment},
    stringstyle=\color{vscodeString},
    numberstyle=\tiny\color{vscodeForeground},
    identifierstyle=\color{vscodeForeground},
    %
    % Keywords (class 1)
    keywords={auto, break, case, catch, class, const, constexpr, continue,
        default, delete, do, else, enum, explicit, extern, false, final,
        for, friend, goto, if, inline, mutable, namespace, new, noexcept,
        nullptr, operator, override, private, protected, public, register,
        return, sizeof, static, static_assert, struct, switch, template,
        this, throw, true, try, typedef, typeid, typename, using, virtual,
        volatile, while},
    %
    % Types (class 2) – highlighted with emph
    emph={bool, char, char16_t, char32_t, double, float, int, long, short,
          signed, unsigned, void, wchar_t, size_t, ptrdiff_t},
    emphstyle={\color{vscodeType}\bfseries},
    %
    % Preprocessor (class 3)
    emph={[2]define, elif, else, endif, error, if, ifdef, ifndef, include,
          line, pragma, undef},
    emphstyle={[2]\color{vscodePreprocessor}\bfseries},
    %
    % Standard library functions (class 4)
    emph={[3]printf, scanf, fprintf, fscanf, sprintf, sscanf, getchar, putchar,
          gets, puts, fopen, fclose, fread, fwrite, malloc, calloc, realloc,
          free, exit, abort, assert, system, getenv, qsort, bsearch, memcpy,
          memset, memcmp, strlen, strcpy, strncpy, strcat, strncat, strcmp,
          strncmp, strstr, strtok, atoi, atof, atol, rand, srand, time, clock,
          pow, sqrt, exp, log, log10, sin, cos, tan, asin, acos, atan, atan2,
          fabs, floor, ceil, fmod, feof, ferror, clearerr, remove, rename,
          tmpfile, tmpnam, setbuf, setvbuf, rewind, fgetc, fputc, ungetc,
          fgets, fputs},
    emphstyle={[3]\color{vscodeFunction}\bfseries},
    %
    % C++ extra keywords (add to class 1)
    morekeywords={alignas, alignof, and, and_eq, bitand, bitor, compl,
        concept, const_cast, decltype, dynamic_cast, module, not, not_eq,
        or, or_eq, reinterpret_cast, requires, static_cast, xor, xor_eq},
    %
    % STL types and functions (add to class 4 via emph again)
    emph={[4]std, vector, map, set, string, array, deque, list, queue, stack,
          pair, tuple, unique_ptr, shared_ptr, weak_ptr, make_shared, make_unique,
          move, forward, function, bind, cout, cin, endl, cerr, clog,
          ifstream, ofstream, fstream, stringstream, istringstream, ostringstream,
          sort, find, transform, accumulate, for_each, copy, reverse, unique,
          remove, replace, swap, min, max, lower_bound, upper_bound, binary_search},
    emphstyle={[4]\color{vscodeFunction}\bfseries},
    %
    frame=single,
    rulecolor=\color{vscodeForeground},
    breaklines=true,
    numbers=left,
    numbersep=5pt,
    showstringspaces=false,
    tabsize=4,
    captionpos=b,
}

% ============================================
% SET DEFAULTS
% ============================================
\lstset{style=vscodepython, language=Python}


%% EMPHASIS BOX

\newtcolorbox{emphasis}{
  colback=white,        % White background
  colframe=black,       % Black border
  boxrule=0.5pt,        % Border thickness
  arc=0pt,              % Sharp corners (0pt = no rounding)
  left=6pt,             % Left padding
  right=6pt,            % Right padding
  top=6pt,              % Top padding
  bottom=6pt,           % Bottom padding
  boxsep=0pt,           % No extra separation
  before upper={
    \setlength{\abovedisplayskip}{0pt}
    \setlength{\abovedisplayshortskip}{0pt}
    \setlength{\belowdisplayskip}{0pt}
    \setlength{\belowdisplayshortskip}{0pt}
  }
}

%%%%%%%%%%%%%%%%%%%%%%%%%%%%%%%%%%%%%%%%%%%%
%             DOCUMENT SETTINGS
%%%%%%%%%%%%%%%%%%%%%%%%%%%%%%%%%%%%%%%%%%%%

\setlength{\parindent}{2em}
\setlength{\parskip}{5pt}
\newlength{\originalparindent}
\newlength{\originalparskip}
\AtBeginDocument{%
    \setlength{\originalparindent}{\parindent}%
    \setlength{\originalparskip}{\parskip}%
}
\numberwithin{equation}{section}

% SECTION TITLE FONT SIZES (move here from preset)
\titleformat*{\section}{\LARGE\bfseries}
\titleformat*{\subsection}{\Large\bfseries}
\titleformat*{\subsubsection}{\large\bfseries}
\titleformat*{\subsubsubsection}{\large\bfseries}

% CUSTOM SECTION COMMANDS
\newcommand{\module}[1]{%
    \clearpage
    \refstepcounter{section}
    \vspace*{1cm}
    {\LARGE\bfseries\noindent Module \thesection\par}
    \vspace{0.5cm}
    {\noindent\Huge\bfseries #1\par}
    \vspace{2cm}
    \addcontentsline{toc}{section}{Module \thesection : #1}
}

\newcommand{\unitsec}[1]{%
    \clearpage
    \refstepcounter{section}
    \vspace*{1cm}
    {\LARGE\bfseries\noindent Unit \thesection\par}
    \vspace{0.5cm}
    {\noindent\huge\bfseries #1\par}
    \vspace{1.5cm}
    \addcontentsline{toc}{section}{Unit \thesection : #1}
}

\newcommand{\chaptersec}[2][\relax]{%
    \clearpage
    \refstepcounter{section}
    \vspace*{0.5cm}
    {\LARGE\bfseries\noindent Chapter \thesection\par}
    \vspace{0.5cm}
    {\noindent\huge\bfseries
        \ifx\relax#1\relax
            #2%
        \else
            #1%
        \fi
        \par
    }
    \vspace{1.5cm}
    \addcontentsline{toc}{section}{Chapter \thesection : #2}
}

%%% PROBLEM ENVIRONMENT (was in preset.tex, now here)
\newenvironment{problem}{%
    \hrule
    \vspace{5pt}
}{%
    \vspace{10pt}\hrule\vspace{10pt}
}

\newenvironment{solution_pset}{%
    \noindent\textbf{\Large Solution}\\[0.5cm]
}{%
    \newpage
}


%%% DEFINITION ENVIRONMENT
% Definition counter linked to section
\newcounter{definition}[section]
\renewcommand{\thedefinition}{\thesection.\arabic{definition}}

\newtcolorbox{definition}[2][]{%
    colback=white,
    colframe=black,
    boxrule=0.5pt,
    arc=0pt,
    left=6pt,
    right=6pt,
    top=6pt,
    bottom=6pt,
    fonttitle=\bfseries,
    title={%
        \stepcounter{definition}% Increment counter
        Definition \thedefinition{} : #2% Then display
    },
    #1
}

%%% EXAMPLE ENVIRONMENT
% Example counter linked to section
\newcounter{example}[section]
\renewcommand{\theexample}{\thesection.\arabic{example}}

\newtcolorbox{example}[2][]{%
    breakable,
    colback=white,
    colframe=blue,
    boxrule=0.5pt,
    arc=4pt,
    left=6pt,
    right=6pt,
    top=6pt,
    bottom=6pt,
    fonttitle=\bfseries,
    parbox=false,  % Don't use parbox mode
    before upper={%
        \setlength{\parindent}{\originalparindent}%
        \setlength{\parskip}{\originalparskip}%
        \noindent%
    },
    title={%
        \refstepcounter{example}%
        Example \theexample{} : #2%
    },
    #1
}

%%% THEOREM ENVIRONMENT
% Theorem counter linked to section
\newcounter{theorem}[section]
\renewcommand{\thetheorem}{\thesection.\arabic{theorem}}
\definecolor{darkred}{RGB}{139, 0, 0} 
\newenvironment{theorem}[2][]{%
    \refstepcounter{theorem}%
    \begin{tcolorbox}[
        colback=white,
        colframe=darkred,
        coltitle=white,
        colbacktitle=darkred,
        boxrule=0.5pt,
        arc=0pt,
        left=6pt,
        right=6pt,
        top=6pt,
        bottom=6pt,
        fonttitle=\bfseries,
        title={Theorem \thetheorem{} : #2},
        #1
    ]
}{%
    \end{tcolorbox}%
}

%%% NOTE ENVIRONMENT
\definecolor{darkgray}{RGB}{100, 100, 100}

\newtcolorbox{note}{%
    colback=white,
    colframe=darkgray,
    boxrule=0.5pt,
    arc=0pt,
    left=6pt,
    right=6pt,
    top=6pt,
    bottom=6pt,
    before upper={%
        \noindent{\textit{Note : }}\par\vspace{1em}%
    }
}

%%% REMARK ENVIRONMENT
\newtcolorbox{remark}{%
    colback=white,
    colframe=darkgray,
    boxrule=0.5pt,
    arc=0pt,
    left=6pt,
    right=6pt,
    top=6pt,
    bottom=6pt,
    before upper={%
        \noindent\textbf{\textit{Remark : }}\par\vspace{1em}%
    }
}

\usepackage{bbm}
\begin{document}
	\begin{titlepage}
    \vspace*{2.5in}
    \hrule
    \vspace{3pt}
    \hrule
    \vspace{0.25in}
    \begin{center}
		{
			\Large 
			\textsc{
				University of the Philippines Diliman
			}
			\\[4pt]
			College of Science
			\\[4pt]
			National Institute of Physics
		}
		\\[0.5cm]
		{
			\Huge
			\textbf{
				Physics 180 THV-1
				\\[5pt]
				Problem Set 8-9
			}
		}
		\\[0.75cm]
		{
			\Large 
			\textsc{
				Vercil Juan
			}
		}\\[5pt] 
		{
			\large 
			IV - BS Physics
		}
    \end{center}
    \vspace{0.25in}
    \hrule
    \vspace{3pt}
    \hrule 
    \vfill
    \begin{center}
    {\large 
		September 17, 2026
	}
    \end{center}
\end{titlepage}

	% ===========================
	% PROBLEM SET 8
	% ===========================
	\stepcounter{section}
	\section*{
		Problem 8.1
	}
	\begin{problem}
		Particle $A$, at rest decays into particles $B$ and $C$:
		\begin{equation}
			\label{eq:decay}
			A \to B+C
		\end{equation}
		What are the energies $E_B$ and $E_C$ of the outgoing particles $B$ and $C$ in 
		terms of the various masses? What are the magnitudes of the outgoing 
		momenta $\abs{p_B}$ and $\abs{p_C}$ (also in terms of the masses)? As a test to see if 
		your answer is correct, $\abs{p_B}$
		should go to zero when $m_A = m_B + m_C$, 
		and becomes imaginary when $m_A < (m_B + m_C)$. Explain why these 
		behaviors happen.
	\end{problem}
	\begin{solution_pset}
		\nocite{flores2026}
		In the rest frame of $A$, four-momentum conservation gives
		\[
			p^\mu_A = p^\mu_B + p_C^\mu
		\]
		with
		\[
			p^\mu_A = (m_A,\vb{0})
		\]
		in natural units $c=1$. Since the outgoing momenta must be equal and opposite,
		\[
			p_B = -p_C, \qquad \abs{p_B}=\abs{p_C}\equiv p
		\]
		The energies are therefore
		\[
			E_B = \sqrt{p^2 + m_B^2},\qquad E_C = \sqrt{p^2 + m_C^2}
		\]
		and energy conservation requires
		\[
			E_B + E_C = m_A
		\]
		Solving for this gives
		\[
			\boxed{
				E_B = \frac{m_A^2 + m_B^2 - m^2_C}{2m_A} 
			}
		\]
		and
		\[
			\boxed{
				E_C = \frac{m_A^2 + m_C^2 -m_B^2}{2m_A} 
			}
		\]
		The common magnitude of the outgoing momentum is
		\[
			p = \sqrt{E_B^2 - m_B^2}
		\]
		Substituting the expression for $E_B$, we have
		\[
			p = \frac{1}{2m_A}\sqrt{
				\pqty{
					m_A^2 - (m_B + m_C)^2
				}
				\pqty{
					m_A^2 - (m_B - m_C)^2
				}
			} 
		\]
		Thus
		\[
			\boxed{
				\abs{p_B} = \abs{p_C} = 
				\frac{
					\sqrt{
						\pqty{
							m_A^2 - (m_B + m_C)^2
						}
						\pqty{
							m_A^2 - (m_B - m_C)^2
						}
					}
				}{
					2M_A
				} 
			}
		\]
		The threshold behaviour, i.e., if $m_A = m_B + m_C$, then
		\[
			m_A^2 - (m_B^2 + m_C^2)^2 = 0
		\]
		so
		\[
			p = 0
		\]
		Physically, this is the threshold for decay.
		The entire rest energy of $A$ is just enough
		to provide the rest masses of $B$ and $C$. There is 
		no energy left for kinetic energy. 

		If $m_A < m_B + m_C$, then
		\[
			m_A^2 - (m_B^2 + m_C)^2 < 0
		\]
		while the other factor remains positive for positive masses
		in the relevant region. Consequently
		\[
			p^2 < 0
		\]
		and $p$ becomes imaginariy. This means there is no physical kinematic solution.
		A does not have enough invariant mass to produce $B+C$, therefore the decay 
		is forbidden.
	\end{solution_pset}
	\newpage

	% ===========================
	% PROBLEM SET 9
	% ===========================
	\stepcounter{section}
	\section*{
		Problem 9.1
	}
	\begin{problem}
		Construct another explicit representation of the $\alpha$ and $\beta$ matrices
		different from the Dirac-Pauli representation and show that they satisfy the anti-commutation
		relations 
		\begin{align}
			\alpha_x^2 =
			\alpha_y^2 =
			\alpha_z^2 =
			\beta^2 
			&= \mathbbm{1}\\
			\alpha_j \beta + \beta \alpha_j 
			&= 0\\
			\alpha_j \alpha_k 
			+
			\alpha_k \alpha_j 
			&= 0\qquad (j\neq k) 
		\end{align}
		How are your new representations related to the Dirac-Pauli
		representation? Does the specific form of the Dirac-Pauli representation
		offer any advantages (if any) to the one you've constructed and versus
		any other representations in general?
	\end{problem}
	\begin{solution_pset}
		The relations require the matrices to have even dimension. Since
		\[
			\beta^2=\mathbbm{1},
		\]
		the eigenvalues of $\beta$ can only be $+1$ and $-1$. Let $V_+$ and $V_-$ denote the corresponding eigenspaces. If $\psi_+\in V_+$, then
		\[
			\beta\psi_+=\psi_+.
		\]
		Using
		\[
			\alpha_j\beta+\beta\alpha_j=0,
		\]
		we obtain
		\[
			\beta(\alpha_j\psi_+)
			=-\alpha_j\beta\psi_+
			=-\alpha_j\psi_+.
		\]
		Thus, $\alpha_j$ maps $V_+$ into $V_-$. Similarly, it maps $V_-$ into $V_+$. Since
		\[
			\alpha_j^2=\mathbbm{1},
		\]
		each $\alpha_j$ is invertible, so these mappings must be one-to-one in both directions. Therefore,
		\[
			\dim V_+=\dim V_-.
		\]
		If each eigenspace has dimension $n$, then the total dimension is
		\[
			N=\dim V_++\dim V_-=2n.
		\]
		Hence the matrices can only have even dimension:
		\[
			N=2,4,6,\ldots.
		\]
		In particular, a $3\times3$ representation is impossible.

		Although $2\times2$ is allowed by this argument, it is not large enough to represent all four matrices $\alpha_x,\alpha_y,\alpha_z,\beta$. The three Pauli matrices already provide three mutually anticommuting $2\times2$ matrices,
		\[
			\sigma_x=
			\begin{pmatrix}
				0&1\\
				1&0
			\end{pmatrix},
			\qquad
			\sigma_y=
			\begin{pmatrix}
				0&-i\\
				i&0
			\end{pmatrix},
			\qquad
			\sigma_z=
			\begin{pmatrix}
				1&0\\
				0&-1
			\end{pmatrix},
		\]
		which satisfy
		\[
			\sigma_j^2=\mathbbm{1},
			\qquad
			\sigma_j\sigma_k+\sigma_k\sigma_j=0
			\qquad (j\neq k).
		\]
		However, there is no fourth nonzero $2\times2$ matrix that anticommutes with all three Pauli matrices while also satisfying $M^2=\mathbbm{1}$. Thus $2\times2$ is insufficient for the full set of four matrices. The next possible dimension is therefore $4\times4$, which is the smallest nontrivial representation of the Dirac algebra.

		One representation different from the Dirac--Pauli representation is the Weyl, or chiral, representation\cite{Weyl1929}. Define
		\[
			\alpha_j=
			\begin{pmatrix}
				\sigma_j&0\\
				0&-\sigma_j
			\end{pmatrix},
			\qquad
			\beta=
			\begin{pmatrix}
				0&\mathbbm{1}\\
				\mathbbm{1}&0
			\end{pmatrix}.
		\]
		Explicitly,
		\[
			\alpha_x=
			\begin{pmatrix}
				0&1&0&0\\
				1&0&0&0\\
				0&0&0&-1\\
				0&0&-1&0
			\end{pmatrix},
			\qquad
			\alpha_y=
			\begin{pmatrix}
				0&-i&0&0\\
				i&0&0&0\\
				0&0&0&i\\
				0&0&-i&0
			\end{pmatrix},
		\]
		\[
			\alpha_z=
			\begin{pmatrix}
				1&0&0&0\\
				0&-1&0&0\\
				0&0&-1&0\\
				0&0&0&1
			\end{pmatrix},
			\qquad
			\beta=
			\begin{pmatrix}
				0&0&1&0\\
				0&0&0&1\\
				1&0&0&0\\
				0&1&0&0
			\end{pmatrix}.
		\]

		Since the Pauli matrices satisfy $\sigma_j^2=\mathbbm{1}$,
		\[
			\alpha_j^2
			=
			\begin{pmatrix}
				\sigma_j^2&0\\
				0&\sigma_j^2
			\end{pmatrix}
			=
			\begin{pmatrix}
				\mathbbm{1}&0\\
				0&\mathbbm{1}
			\end{pmatrix}
			=\mathbbm{1}_4.
		\]
		Furthermore,
		\[
			\beta^2
			=
			\begin{pmatrix}
				0&\mathbbm{1}\\
				\mathbbm{1}&0
			\end{pmatrix}
			\begin{pmatrix}
				0&\mathbbm{1}\\
				\mathbbm{1}&0
			\end{pmatrix}
			=
			\begin{pmatrix}
				\mathbbm{1}&0\\
				0&\mathbbm{1}
			\end{pmatrix}
			=\mathbbm{1}_4.
		\]
		Therefore,
		\[
			\alpha_x^2=\alpha_y^2=\alpha_z^2=\beta^2=\mathbbm{1}_4.
		\]

		For the anticommutation relation between $\alpha_j$ and $\beta$,
		\[
			\alpha_j\beta
			=
			\begin{pmatrix}
				\sigma_j&0\\
				0&-\sigma_j
			\end{pmatrix}
			\begin{pmatrix}
				0&\mathbbm{1}\\
				\mathbbm{1}&0
			\end{pmatrix}
			=
			\begin{pmatrix}
				0&\sigma_j\\
				-\sigma_j&0
			\end{pmatrix},
		\]
		while
		\[
			\beta\alpha_j
			=
			\begin{pmatrix}
				0&\mathbbm{1}\\
				\mathbbm{1}&0
			\end{pmatrix}
			\begin{pmatrix}
				\sigma_j&0\\
				0&-\sigma_j
			\end{pmatrix}
			=
			\begin{pmatrix}
				0&-\sigma_j\\
				\sigma_j&0
			\end{pmatrix}.
		\]
		Hence,
		\[
			\alpha_j\beta+\beta\alpha_j
			=
			\begin{pmatrix}
				0&0\\
				0&0
			\end{pmatrix}
			=0.
		\]

		For $j\neq k$,
		\[
			\alpha_j\alpha_k
			=
			\begin{pmatrix}
				\sigma_j&0\\
				0&-\sigma_j
			\end{pmatrix}
			\begin{pmatrix}
				\sigma_k&0\\
				0&-\sigma_k
			\end{pmatrix}
			=
			\begin{pmatrix}
				\sigma_j\sigma_k&0\\
				0&\sigma_j\sigma_k
			\end{pmatrix},
		\]
		and
		\[
			\alpha_k\alpha_j
			=
			\begin{pmatrix}
				\sigma_k\sigma_j&0\\
				0&\sigma_k\sigma_j
			\end{pmatrix}.
		\]
		Since
		\[
			\sigma_j\sigma_k+\sigma_k\sigma_j=0
			\qquad (j\neq k),
		\]
		we have
		\[
			\alpha_j\alpha_k+\alpha_k\alpha_j
			=
			\begin{pmatrix}
				\sigma_j\sigma_k+\sigma_k\sigma_j&0\\
				0&\sigma_j\sigma_k+\sigma_k\sigma_j
			\end{pmatrix}
			=0.
		\]
		Thus the Weyl representation satisfies all of the required relations.

		The Dirac--Pauli representation is
		\[
			\alpha_j^{(D)}
			=
			\begin{pmatrix}
				0&\sigma_j\\
				\sigma_j&0
			\end{pmatrix},
			\qquad
			\beta^{(D)}
			=
			\begin{pmatrix}
				\mathbbm{1}&0\\
				0&-\mathbbm{1}
			\end{pmatrix}.
		\]
		The Weyl and Dirac--Pauli representations are related by a unitary change of basis. In particular, with
		\[
			U=
			\frac{1}{\sqrt{2}}
			\begin{pmatrix}
				\mathbbm{1}&\mathbbm{1}\\
				-\mathbbm{1}&\mathbbm{1}
			\end{pmatrix},
		\]
		the Weyl matrices are obtained from the Dirac--Pauli matrices through
		\[
			\alpha_j^{(W)}
			=
			U\alpha_j^{(D)}U^\dagger,
			\qquad
			\beta^{(W)}
			=
			U\beta^{(D)}U^\dagger.
		\]
		Thus, the two representations contain the same physical information and differ only in the choice of basis used to represent the Dirac spinors.

		The particular form of the Dirac--Pauli representation can nevertheless be advantageous for certain calculations. In the Dirac--Pauli representation,
		\[
			\beta^{(D)}
			=
			\begin{pmatrix}
				\mathbbm{1}&0\\
				0&-\mathbbm{1}
			\end{pmatrix},
		\]
		so the upper and lower two-component parts of a Dirac spinor are separated naturally. This makes the nonrelativistic limit particularly convenient, since the upper components of a positive-energy solution are the dominant components.

		In contrast, the Weyl representation has
		\[
			\alpha_j^{(W)}
			=
			\begin{pmatrix}
				\sigma_j&0\\
				0&-\sigma_j
			\end{pmatrix},
		\]
		which makes the two chiral components explicit. It is consequently particularly convenient for calculations involving chirality and massless fermions.

		Therefore, no representation is physically preferred over the others. The different representations are related by changes of basis and give the same physical predictions. Their advantages are instead computational: the Dirac--Pauli representation is convenient for the nonrelativistic limit, while the Weyl representation is convenient for chirality and massless particles.
	\end{solution_pset}
	\newpage

	% ===========================
	\stepcounter{section}
	\bibliographystyle{unsrturl} 
	\bibliography{references} 
	\stepcounter{section}
	% ===========================
	\section*{Appendices}
		\label{appendix}
	
		\subsection*{Code} 
		No code was used for this problem set
	\vfill
	\begin{center}
		\textit{Thank you Mr. Kong}\\
		\vspace{1cm}
		{\large
		\textsc{Physica Vincit Omnia}
		}
	\end{center}
\end{document}
